\documentclass[11pt]{article}

\usepackage[a4paper,margin=2.4cm]{geometry}
\usepackage{amsmath}
\usepackage{booktabs}
\usepackage{enumitem}
\usepackage[hidelinks]{hyperref}
\usepackage{siunitx}

\newcommand{\file}[1]{\texttt{#1}}

\title{Shared plan for data-driven emitter generation and\\
weighted list-mode propagation}
\author{J.J. G\'omez-Cadenas}
\date{\today}

\begin{document}
\maketitle

\begin{abstract}
This note defines the interface between the Geant4 proton and positron
calculation in \file{ptcryspg4} and the PET detector calculation in
\file{PTCryspMC.jl}.  Three source-generation modes are retained: native
Geant4 residual production for historical compatibility, an unweighted source
generated from the nominal experimental cross-section curves for routine
data-driven studies, and a shared weighted source bank for propagation of the
full cross-section ensemble.  The first two modes produce ordinary emitter
lists.  The third transports each sampled source and detector event once and
assigns curve-dependent weights afterwards.  This document records the common
method, data contracts and division of responsibilities between the two
repositories.
\end{abstract}

\section{Purpose and relationship to the other notes}

The experimental data selection, curation, fitted excitation functions and
1000 correlated replicas are documented in \file{docs/xsection\_fit.tex}.
The proton-exposure estimator and its Geant4 implementation are specified in
\file{docs/xsections\_plan.tex}.  The present document is the cross-repository
plan: it describes how a production event is carried from the Geant4
calculation to detector-level list-mode lines of response (LORs), and how the
nominal and replica-dependent event populations are represented.

Implementation details local to \file{PTCryspMC.jl} will be recorded in a
separate developer note in that repository.  Its generated
\file{docs/SCHEMA.md} remains the authoritative description of the implemented
singles and LOR formats.

All implementation changes are made on dedicated feature branches.  The G4
work uses \file{xsection-source-bank}; the detector-side work will use
\file{xsection-weighted-lors}.  The existing production applications remain
available while the new calculation is validated.

\section{Terminology}

The earlier plan used \emph{proposal distribution} and \emph{proposal source},
which are standard importance-sampling terms but obscure the role of the data
product.  The implementation and user documentation will use the following
terms:

\begin{description}[style=nextline,leftmargin=0pt]
  \item[Shared source bank]
  A reusable collection of candidate emitter events covering the selected
  targets, production channels, proton energies and spatial regions.

  \item[Bank sampling distribution]
  The probability rule used to select the candidate events.  It has support
  wherever the nominal curve or any retained replica can produce an emitter.

  \item[Source-bank mode]
  The Geant4 mode that generates the shared source bank and transports the
  positron associated with each bank event.

  \item[Weighted source event]
  A bank event after assigning its contribution for the nominal curve or one
  cross-section replica.
\end{description}

The symbol $q_{s,c}$ remains useful in equations as the bank sampling
probability for proton segment $s$ and production channel $c$.  The word
``proposal'' is unnecessary in filenames, run modes and user-facing text.

\section{Three source-generation modes}

The three modes are scientifically distinct and mutually exclusive for a
given source calculation.  Table~\ref{tab:modes} summarizes their roles.

\begin{table}[ht]
\centering\small
\caption{Source-generation modes.}
\label{tab:modes}
\begin{tabular}{p{0.18\textwidth}p{0.22\textwidth}p{0.25\textwidth}p{0.25\textwidth}}
\toprule
Mode & Production cross sections & Principal purpose & Detector input \\
\midrule
Native G4 & Geant4 physics-list residual production & Reproduce the
established Stage-A calculation and provide a model comparison & Ordinary
unweighted \file{emitters.csv} plus the existing handoff budget \\
Data-driven nominal & Median experimental excitation functions & Routine
best-estimate production studies & Ordinary unweighted
\file{emitters.csv} plus a nominal data-driven handoff budget \\
Shared weighted bank & Nominal curve and all correlated replicas & Propagate
cross-section uncertainty with common source and detector histories & Shared
source-event table and detector LOR bank with stable event identifiers \\
\bottomrule
\end{tabular}
\end{table}

\subsection{Mode 1: native Geant4}

The current \file{stageA\_transport} application defines the native-G4 mode.
Geant4 selects the inelastic interaction, residual isotope and final state from
the chosen physics list.  Its standard radioactive-decay and positron
transport produce the annihilation point.  The existing
\file{emitters.csv}, \file{run\_meta.csv}, depth-dose and handoff products keep
their present meanings.

This mode provides exact backward compatibility with earlier studies.  It is
also a useful comparison because differences between it and the data-driven
nominal source measure the practical consequence of replacing the Geant4
partial residual-production model with the fitted experimental excitation
functions.  Native-G4 results are identified explicitly in filenames,
metadata, plots and tables.

The existing application remains the implementation of this mode.  The new
cross-section application may read the same versioned beam and geometry
inputs, while agreement is established through dose and depth-dose closure.

\subsection{Mode 2: data-driven nominal}

The nominal mode uses the median fitted experimental curve
$\sigma_c^{\mathrm{nom}}(E)$ for each selected channel.  For proton segment
$s$ and channel $c$, its expected production contribution is
\begin{equation}
 \Delta N_{s,c}^{\mathrm{nom}}
 = a_s n_{t(c),s}\ell_s\sigma_c^{\mathrm{nom}}(E_s),
 \label{eq:nominal-contribution}
\end{equation}
where $a_s$ is the proton statistical weight, $n_{t(c),s}$ is the target-isotope
number density and $\ell_s$ is the step length.

Production sites and channels are sampled in proportion to
Eq.~\eqref{eq:nominal-contribution}.  The associated residual isotope is
created at the sampled production position and its positron is transported to
annihilation.  The resulting per-isotope point pools are unweighted.  Their
physical production budgets are obtained from the folded nominal yields and
the established per-Gy normalization.  The beam-off handoff applies the named
irradiation, delay and acquisition scenario.

This mode is the routine data-driven calculation.  It writes the ordinary core
\file{emitters.csv} interface and an ordinary sampling budget, so the existing
\file{PTCryspMC.jl} API path can consume it without cross-section weights.  The
data-driven provenance identifies the fit version and nominal-curve digest.

The data-driven nominal result can differ substantially from the native-G4
result, particularly where the Geant4 residual model has incomplete support.
The nominal experimental fit defines the best-estimate production source;
the native-G4 result remains a separately labelled comparison.

\subsection{Mode 3: shared weighted source bank}

The uncertainty calculation uses a common bank of sampled proton-segment and
channel combinations.  The bank sampling distribution $q_{s,c}$ covers the
support of the nominal fit and the complete retained replica ensemble.  For
curve $r$, bank event $e=(s,c)$ carries weight
\begin{equation}
 w_e^{(r)}=
 \frac{a_s n_{t(c),s}\ell_s\sigma_c^{(r)}(E_s)}{q_{s,c}}.
 \label{eq:event-weight}
\end{equation}
Common overall sampling and per-Gy factors are recorded in the bank metadata
and applied consistently to every curve.

Each bank event fixes the production channel, isotope, production point and
proton energy.  Geant4 transports its positron once and records the
annihilation point.  The PET calculation then transports the two annihilation
photons through the phantom and scanner once.  The event consequently has one
stored detector response and one weight for every cross-section curve.

The nominal curve is included in this weighted calculation.  Its weighted
source and reconstructed profiles provide the reference for all paired
replica displacements.  Mode 2 remains the convenient unweighted route for
normal studies and provides an independent closure test: the weighted nominal
profile from mode 3 must agree with the unweighted nominal profile from mode 2
within the predefined Monte Carlo tolerances.

The weighted bank carries the replica ensemble through the detector
calculation.  Regenerating an unweighted source for one selected curve
retraces the same bank events and validates the weighted result for that curve
at full detector level.

The weighted method changes the spatial density of LOR contributions.  The
weight depends on channel and proton energy, both of which vary with depth.
The stored LOR endpoints remain fixed while the weighted activity profile and
its distal $R_{50}$ change from one curve to another.

\section{Shared source-event contract}

The source bank requires one stable record per transported positron.  The
record contains at least
\begin{itemize}[nosep]
  \item a stable source-event identifier;
  \item the residual isotope and production-channel identifiers;
  \item the target isotope;
  \item the proton kinetic energy or a reference to its validated energy bin;
  \item the production and annihilation coordinates;
  \item the bank sampling probability or an equivalent precomputed base
        factor;
  \item the proton-exposure run identifier and digest;
  \item the cross-section fit and bank-definition digests.
\end{itemize}

The source-event table stores these values once.  Detector singles and LORs
carry compact source-event identifiers rather than duplicating the complete
record.  Replica weights are evaluated from the source-event table and the
cross-section lookup tables.  A matrix containing one stored weight per LOR
and replica is therefore unnecessary.

The production and annihilation coordinates have different roles.  The
production position defines the activity-production profile and its
$R_{50,\mathrm{prod}}$.  The annihilation position is the source of the two
\SI{511}{keV} photons and defines the position entering the detector
simulation.

\section{PTCryspMC.jl propagation}

The current detector pipeline is
\begin{center}
source events $\longrightarrow$ singles $\longrightarrow$ truth LORs and
random LORs $\longrightarrow$ detector LORs.
\end{center}
The existing event number already propagates through singles and same-decay
truth/scatter LORs.  It can serve as the source-event lookup key when its
mapping to the shared source-event table is preserved.

For a true or scattered LOR $L$ produced by source event $e$,
\begin{equation}
 W_L^{(r)}=w_e^{(r)}.
 \label{eq:true-lor-weight}
\end{equation}
A random LOR combines photons from two distinct source events.  Its schema
therefore requires both event identifiers.  If its photons originate from
$e_1$ and $e_2$, its expected contribution is
\begin{equation}
 W_{L,\mathrm{random}}^{(r)}=
 w_{e_1}^{(r)}w_{e_2}^{(r)}.
 \label{eq:random-lor-weight}
\end{equation}
The existing random builder already has both identifiers while forming the
pair; the extended schema preserves them through detector smearing and
selection.

The random pair list is formed once, at the nominal activity.  Reweighting
this fixed list describes a replica's accidental contribution to first order
in its activity change; the weighted random-rate closure in the acceptance
list measures the residual of this approximation.

The detector code will retain two input paths:
\begin{enumerate}[nosep]
  \item the existing unweighted API source for native-G4 and data-driven
        nominal scenarios; and
  \item an optional shared-bank source that preserves source-event identity
        through the detector list.
\end{enumerate}

The reconstruction stage can consume weighted list-mode events directly.  A
tracked resampling program can alternatively materialize ordinary unweighted
LOR lists for a selected replica subset.  Both routes reuse the same photon
transport, detector smearing and selection realization.  The chosen route and
all weight-table digests form part of the reconstruction provenance.

\section{Beam-off handoff and normalization}

All three modes retain the established division of responsibilities.  Geant4
determines production and annihilation space.  The handoff determines the
expected number of measured decays for a named irradiation and acquisition
scenario.  For isotope $j$, its production contribution is multiplied by the
same build-up, delay and acquisition-window factors already used by
\file{decay\_sampling/budget.py}.

In modes 1 and 2, the handoff writes an ordinary per-isotope sampling budget.
In mode 3, the isotope-dependent handoff factor multiplies each source-event
weight.  The named handoff scenario and its file digest are common to the G4
folding products, detector source and reconstructed LOR analysis.

Per-proton and per-Gy normalization use the proton-exposure run metadata.  The
target dose convention is shared with the current Stage-A budget calculation.
The native-G4, unweighted nominal and weighted nominal outputs are all reported
for the same primary count, target-dose definition and beam-off scenario.

\section{Geant4 work still required}

The experimental curves, replica ensemble, folding implementation, synthetic
exposure tests and forced-final-state Geant4 diagnostic already exist.  The
physical application \file{stageA\_xsection\_ensemble/} remains to be built.
Its required components are:

\begin{enumerate}
  \item An independent Geant4 application with its own CMake target, geometry,
        primary generator, actions, run macros and tests.

  \item A proton-exposure mode scoring primary and secondary proton track
        length by target isotope, energy and position, together with dose,
        depth dose, native-route diagnostics and complete run metadata.

  \item A source-bank mode implementing a documented bank sampling
        distribution, stable event identifiers, channel and energy retention,
        positron transport and the shared source-event table.

  \item A data-driven nominal mode producing unweighted emitter pools and
        nominal handoff budgets from the median experimental curves.

  \item Runtime selection of BIC, INCL++ and Bertini for the independent
        proton-exposure envelope.  The external production curves remain
        common to these runs.

  \item Closure tests connecting per-step exposure, binned folding,
        unweighted nominal sampling and weighted-bank sampling.
\end{enumerate}

\section{Validation and acceptance}

The shared calculation is accepted when the following conditions hold:

\begin{enumerate}
  \item The native-G4 mode reproduces the existing Stage-A output and handoff
        semantics.

  \item The unweighted nominal source reproduces the nominal folded isotope
        yields and production profiles.

  \item The weighted nominal source reproduces the unweighted nominal source
        in isotope yield, production $R_{50}$, annihilation $R_{50}$ and
        reconstructed $R_{50}$ within the frozen statistical tolerances.

  \item Every replica reproduces its direct-fold production profile when the
        shared source-event weights are summed before detector transport.

  \item The event-weighted detector result is invariant, within its tolerance,
        under enlargement of the shared source bank.

  \item True and scattered LORs preserve one source identifier; random LORs
        preserve both source identifiers and satisfy the expected weighted
        random-rate closure.

  \item A schema-version guard prevents mixing native, nominal and shared-bank
        files with incompatible semantics.

  \item All reported calculations, comparisons and plots are generated by
        tracked repository code.
\end{enumerate}

\section{Implementation sequence across repositories}

\begin{enumerate}
  \item Finalize this shared contract and update
        \file{docs/xsections\_plan.tex} on the G4 feature branch.

  \item Implement and validate the physical proton-exposure scorer in the new
        G4 application.

  \item Implement the data-driven nominal source and establish closure with
        the direct nominal fold.

  \item Implement the shared source bank and establish weighted-nominal versus
        unweighted-nominal closure before detector propagation.

  \item Create the detector-side feature branch and document its exact source,
        singles and LOR schema changes.

  \item Add the shared-bank source mode, source-event lookup and the second
        random-event identifier to \file{PTCryspMC.jl}; preserve the existing
        unweighted API path.

  \item Propagate the weighted nominal and a controlled replica subset through
        the full detector and reconstruction chain, then compare weighted and
        resampled list-mode reconstruction where required.

  \item Run the complete replica ensemble through the validated offline
        weighting layer and report the paired reconstructed-edge distribution.
\end{enumerate}

\end{document}
